\chapter{Gegner-Almanach}

% Wiederverwendbare Darstellung fuer Gegnerprofile

\definecolor{hpEnemy}{HTML}{F2EEE6}
\definecolor{hpEnemyLine}{HTML}{8A6A49}
\definecolor{hpEnemyInner}{HTML}{FBF9F4}

% Ein Gegnerprofil bildet einen geschlossenen Almanach-Eintrag.
\newtcolorbox{enemyprofile}[2][]{
  hpbox,
  title={#2},
  colback=hpEnemy,
  colframe=hpEnemyLine,
  boxed title style={colback=hpEnemyLine!22},
  fonttitle=\HPHeadingFont\large,
  #1
}

% Kompakte Teilbereiche innerhalb eines Gegnerprofils.
\newtcolorbox{enemysectionbox}[1]{
  enhanced,
  colback=hpEnemyInner,
  colframe=hpLine,
  boxrule=0.55pt,
  arc=1.5mm,
  left=3mm,
  right=3mm,
  top=2.5mm,
  bottom=2.5mm,
  before skip=0pt,
  after skip=0pt,
  title={#1},
  fonttitle=\HPUIFont\bfseries,
  fontupper=\HPUIFont,
  coltitle=hpBrown,
  attach title to upper={\quad}
}

% Werteband direkt unter dem Gegnernamen.
\newenvironment{enemystats}
  {\begingroup\HPUIFont\renewcommand{\arraystretch}{1.15}\setlength{\tabcolsep}{4pt}%
   \tabularx{\linewidth}{@{}>{\bfseries}l X >{\bfseries}l X >{\bfseries}l X@{}}}
  {\endtabularx\endgroup}

% Zweispaltiger Inhaltsbereich fuer Angriffe, KI und Sonderregeln.
\newenvironment{enemygrid}
  {\begin{tcbraster}[
      raster columns=2,
      raster column skip=4mm,
      raster row skip=3mm,
      raster valign=top
    ]}
  {\end{tcbraster}}

\newcommand{\enemyrole}[1]{\textbf{Rolle:} #1}
\newcommand{\enemydesign}[1]{\textbf{Designziel:} #1}


Der Gegner-Almanach versammelt die spielrelevanten Werte, Angriffe und Verhaltensregeln der Gegner. Jeder Eintrag ist als kompaktes Profil aufgebaut: Das Werteband zeigt die wichtigsten Kennzahlen, darunter stehen Angriffe und Verhalten direkt nebeneinander.

\section{Höhlenspinne}

\begin{enemyprofile}{Höhlenspinne}
  \begin{enemystats}
    Lebenspunkte & 12~\lifeLoss & Schilde & 1~\armor & Rolle & Kontrollgegner \\
  \end{enemystats}

  \smallskip
  Eine lauernde Jägerin, die ihre Beute zunächst mit klebrigen Netzen schwächt und anschließend aus nächster Nähe zubeißt.

  \begin{enemygrid}
    \begin{enemysectionbox}{Angriffe}
      \textbf{Netzspucken} \hfill \rangedAttack

      5~\orangeDmgDie{} Netzwürfel

      Jeder nicht negierte Erfolg erzeugt einen Netzmarker auf dem Ziel.

      \medskip
      \textbf{Biss} \hfill \meleeAttack

      3~\redDmgDie{} Schadenswürfel
    \end{enemysectionbox}

    \begin{enemysectionbox}{Verhalten}
      Solange mindestens ein Held keinen Netzmarker besitzt, setzt die Höhlenspinne \textbf{Netzspucken} ein.

      Besitzen alle Helden mindestens einen Netzmarker, führt sie einen \textbf{Biss} aus.
    \end{enemysectionbox}

    \begin{enemysectionbox}{Netzmarker}
      Jeder Netzmarker verhindert den nächsten Angriffserfolg oder Ausweicherfolg des betroffenen Helden.

      Anschließend wird der Marker entfernt.
    \end{enemysectionbox}

    \begin{enemysectionbox}{Taktische Rolle}
      \enemyrole{Kontrolle und Vorbereitung}

      \medskip
      \enemydesign{Die Höhlenspinne verursacht zunächst wenig direkten Schaden. Ihre Netze schwächen jedoch gezielt die nächsten Angriffe oder Ausweichversuche und erhöhen so den Druck nachfolgender Aktionen.}
    \end{enemysectionbox}
  \end{enemygrid}
\end{enemyprofile}

\section{Baumdrache}

\begin{enemyprofile}{Baumdrache}
  \begin{enemystats}
    Lebenspunkte & 50~\lifeLoss & Schilde & 3~\armor & Rolle & Bossgegner \\
  \end{enemystats}

  \smallskip
  Ein mächtiger Gegner mit einer vorhersehbaren, aber abwechslungsreichen Angriffsfolge. Wer seine Muster erkennt, kann sich gezielt auf den nächsten Schlag vorbereiten.

  \begin{enemygrid}
    \begin{enemysectionbox}{Angriffe}
      \textbf{Feueratem} \hfill \rangedAttack

      5~\orangeDmgDie{} gegen alle Helden

      \medskip
      \textbf{Biss} \hfill \meleeAttack

      3~\redDmgDie{}

      \medskip
      \textbf{Schwanzschlag} \hfill \meleeAttack

      3~\orangeDmgDie{}; das Ziel wird weggeschleudert

      \medskip
      \textbf{Krallenhieb} \hfill \meleeAttack

      4~\orangeDmgDie{}
    \end{enemysectionbox}

    \begin{enemysectionbox}{Angriffsfolge}
      Der Baumdrache führt pro Runde genau eine Aktion aus:

      \begin{enumerate}[leftmargin=1.5em,itemsep=0.15\baselineskip]
        \item Feueratem
        \item Biss
        \item Schwanzschlag
        \item Krallenhieb
        \item Biss
        \item Schwanzschlag
        \item Krallenhieb
      \end{enumerate}

      Danach beginnt die Folge erneut.
    \end{enemysectionbox}

    \begin{enemysectionbox}{Prolog-Regel}
      Im Prolog endet der Kampf nicht mit dem Tod des Baumdrachen.

      Sobald er insgesamt \textbf{30 Schaden} erlitten hat, flieht er.
    \end{enemysectionbox}

    \begin{enemysectionbox}{Taktische Rolle}
      \enemyrole{Boss mit lernbarer Rotation}

      \medskip
      \enemydesign{Der Baumdrache verbindet Flächenschaden, hohen Einzelschaden und Positionskontrolle. Seine feste Reihenfolge macht die Bedrohung planbar und belohnt vorausschauendes Handeln.}
    \end{enemysectionbox}
  \end{enemygrid}
\end{enemyprofile}
